\chapter{Aplicație: recunoașterea muzicii}
\label{ch:shazam}

Recunoașterea unei piese din câteva secunde de audio, înregistrate cu un
microfon de telefon, peste zgomot de fundal și voci, pare să ceară o
comparație directă cu milioane de piese dintr-o bază de date. Algoritmul
din spatele Shazam~\cite{wang2003} evită exact asta, comparând amprente
scurte și discrete în loc de audio brut.

\section{Cum funcționează amprenta}

Primul pas e o spectrogramă a piesei, exact tipul de reprezentare
timp-frecvență din Capitolul~\ref{ch:wavelets}. Din ea se păstrează doar
vârfurile locale de energie, cele mai puternice puncte din fiecare regiune
de timp-frecvență, alese cu un criteriu de densitate care le distribuie
uniform pe toată durata piesei, cele mai tari câteva din fiecare fereastră,
nu doar cele mai tari câteva sute din tot semnalul.
Rezultatul e o „constelație” de puncte $(t, f)$, suficient de rară ca să
încapă compact, suficient de densă ca să supraviețuiască zgomotului.

Un singur punct din constelație nu spune mare lucru, miile de piese din
bază conțin, fiecare, vârfuri la frecvențe apropiate. Amprenta reală vine
din \emph{perechi}: fiecare punct devine un „punct-ancoră”, legat de câteva
puncte dintr-o „zonă țintă” apropiată în timp și frecvență. Fiecare pereche
(ancoră, țintă) generează un hash format din trei numere, frecvența
ancorei, frecvența țintei și diferența de timp dintre ele,
$(f_1, f_2, \Delta t)$, stocat împreună cu timpul absolut la care apare în
piesă. O piesă de câteva minute produce astfel zeci de mii de hash-uri, nu
un singur „cod” global.

Căutarea unui fragment înregistrat repetă exact procesul, spectrogramă,
constelație, hash-uri, și caută fiecare hash în baza de date. Zgomotul de
fundal adaugă vârfuri false și poate ascunde unele reale, dar tot lasă
destule perechi curate ca zeci de hash-uri să se potrivească exact cu o
singură piesă din bază, iar dacă hash-urile potrivite cad la o diferență de
timp \emph{constantă} față de timpul lor din piesa originală (fragmentul
înregistrat e, de exemplu, exact piesa originală începând din secunda 47),
potrivirea e confirmată statistic, nu doar prin numărul de hash-uri comune.
Pentru că amprenta e locală, fiecare hash descrie doar o mică vecinătate
timp-frecvență, metoda tolerează inclusiv piese suprapuse: hash-urile unei
piese ascultate peste alta rămân, în bună parte, neschimbate, deci ambele
pot fi recunoscute din același fragment.

\section{Demo}

Demo-ul din tur (Figura~\ref{fig:shazam-demo}) arată, mai simplu, punctul de
plecare al ideii: trei acorduri redate pe rând interferează în domeniul
timp, iar semnalul mixt arată vizual de ce identificarea nu se poate face
direct pe formă de undă, notele individuale nu mai sunt separabile cu
ochiul. Textul de pe ecran leagă figura de mecanismul real: vârfurile dintr-o
spectrogramă devin perechi $(f_1, f_2, \Delta t)$ căutate într-o bază de
date, exact procesul descris mai sus.

\begin{figure}[htbp]
\centering
\includegraphics[width=\textwidth]{39-applications-shazam.png}
\caption{Trei acorduri interferând în domeniul timp. Textul explică
  amprentele $(f_1, f_2, \Delta t)$ pe care le caută algoritmul real.}
\label{fig:shazam-demo}
\end{figure}
